\documentclass[00-main.tex]{subfiles}

\begin{document}

\section{Conclusions}
Standard DDM fitting methods treat trial sequences as independent and identically distributed, an assumption routinely violated in naturalistic behavioral experiments where decisions are richly structured in time \citep{urai_structure_2025, vincent_recurrent_2017}. This paper introduces a pseudo-MLE method that provides analytically tractable uncertainty quantification for DDM parameters while remaining valid under both temporal dependence between trials and model misspecification.

\subsubsection*{Theoretical Contributions}
We establish that the pseudo-MLE converges to its pseudo-true coefficients and is asymptotically normal under strict stationarity and strong mixing. The key insight is that stationarity applies to the joint process of covariates and outcomes, not to the decision outcomes alone. By explicitly modeling structured non-stationarities --- such as session-level trends, circadian rhythms, and learning effects --- valid point estimates and uncertainties are recovered even in non-stationary examples, substantially broadening the scope of settings to which the method applies. The resulting sandwich-form covariance estimators \citep{newey_simple_1987, hansen_probability_2022} correct for the failure of the information matrix equality under possible misspecification, providing valid standard errors and confidence intervals even when the working model is an approximation of the true data-generating process.

\subsubsection*{Simulation Results}
Monte Carlo simulations confirm that under correct specification, with either constant or covariate-driven non-stationary parameter, point estimates are unbiased and covariances are accurately recovered, consistent with prior DDM parameter recovery studies \citep{boehm_estimating_2018, matzke_psychological_2009}. Under misspecification induced by unmodeled trial-to-trial parameter variability, non-robust covariance estimators substantially underestimate standard errors and correlations, while robust estimators largely mitigate these biases. Robust covariance estimation is therefore advisable whenever misspecification cannot be ruled out, which \citet{urai_structure_2025} notes is the norm in applied DDM research. More broadly, explicitly modeling sources of parameter variability via covariates is preferable to absorbing variability into unstructured noise parameters, which are difficult to estimate due to unidentifiability and carry no temporally structured information \citep{boehm_estimating_2018, lerche_model_2016}.

\subsubsection*{Data Analysis Results}
Applied to empirical data from a single rat performing a visual motion discrimination task \citep{reinagel_speed_2013}, the method reveals systematic, multi-timescale structure in DDM parameters invisible to standard analyses that assume temporal independence. At the between-day timescale, increasing decision boundaries and drift rate coefficients across the first 35 days track a practice effect consistent with gradual attentional improvement \citep{ashwood_mice_2022, mohammadi_identifying_2025}. Within a day, decision boundaries and non-decision times increase monotonically as the rat becomes satiated, while markedly lower boundaries and faster responses in the first 20 trials reflect extreme thirst at session onset. These findings illustrate that even under constant experimental conditions, DDM parameters vary substantially across multiple timescales --- changes that would be misattributed as noise by methods ignoring trial order. Persistent autocorrelations across days and within sessions confirm that sequential dependencies are present between trials and must be accounted for in inference, consistent with findings across rodent and primate decision making \citep{urai_choice_2019, busse_detection_2011, danskin_exponential_2023}.

\subsubsection*{Limitations}
Misspecification-robust covariances protect against inferential errors when the working model is incomplete, but do not substitute for explicitly modeling important sources of variability, such as history-dependent biases that drive sequential dependencies \citep{nguyen_optimizing_2019, urai_choice_2019} or species-specific forms of trial-to-trial variability \citep{nguyen_different_2022, shevinsky_interaction_2019}. The method also requires that, after covariate adjustment, the joint process is stationary and mixing. There are conditions that are difficult to verify in practice, though our simulations demonstrate that realistic non-stationarities do not strongly distort point or uncertainty estimates even when unmodeled. Finally, the present application demonstrates the types of temporal structure that can be uncovered in decision-making data; generalizing these specific findings requires larger-scale analysis across individuals and species.

\subsubsection*{Future Directions}
This work establishes a methodological foundation for studying dynamic decision processes in naturalistic settings. An important next step is to extend the framework to large-scale, multi-subject datasets to characterize how decision-making adapts across trials, days, and months within and across species \citep{asadpour_rodents_2024, danskin_exponential_2023}. Combining the pseudo-MLE approach with latent-state models that explicitly represent strategy switching, such as hidden Markov models \citep{ashwood_mice_2022, mohammadi_identifying_2025}, could jointly capture structured non-stationarity and discrete changes in decision strategies. In conclusion, by relaxing the standard assumptions of correct specification and trial independence, this framework enables principled and computationally efficient inference about dynamic decision processes under the realistic conditions of naturalistic behavioral experiments.
\end{document}